%%%%%%%%%%  scan idx 431  |  folio 418  %%%%%%%%%%
$T_i$ sur $f_i$; notons $u_0:E_0=(A_0)^d\longrightarrow A_0$ l'homomorphisme défini par la suite $(T_1,\ldots,T_d)$, de sorte que $u$ est déduit de $u_0$ par changement de base de $A_0$ à $A$. Comme la suite $(T_1,\ldots,T_d)$ est régulière (au sens classique), l'homomorphisme $u_0$ est régulier, et on est sous les conditions du lemme, appliqué au changement de base $A_0\longrightarrow A$. Or les $J_0^n/J_0^{n+1}$ sont libres de type fini sur $\mathbb Z$, donc on peut appliquer les résultats du a). Par suite $J^n/J^{n+1}\xleftarrow{\sim}J_0^n/J_0^{n+1}\otimes_{A_0}A$, et $J^n/J^{n+1}$ est libre de type fini.

\textbf{Proposition 1.3.} Soient $A$ un anneau, $E$ un $A$-module projectif de type fini, $u:E\longrightarrow A$ un homomorphisme régulier, $J=u(E)$. Alors l'idéal $J$ est quasi-régulier (EGA IV 16.9.7), c'est-à-dire (EGA IV 16.9.8):
\begin{enumerate}[label=\roman*)]
\item $J$ est de type fini.
\item $J/J^2$ est un $A/J$-module projectif de type fini;
\item l'homomorphisme canonique
\end{enumerate}
\begin{equation}\tag{1.3.1}
 \Sym_{A/J}(J/J^2)\longrightarrow \gr_J(A)=\bigoplus_{n\geq0}J^n/J^{n+1},
\end{equation}
est un isomorphisme.

On a vu dans la démonstration de 1.2 b), que $J/J^2$ est projectif de type fini. Reprenons les notations de 1.2 b) (la question étant locale sur $\Spec(A)$, on suppose encore $E$ libre de type fini). Alors l'homomorphisme canonique (1.3.1) se déduit de l'homomorphisme analogue pour $J_0$ par le changement de base $A_0\longrightarrow A$, d'après les relations 1.2. Or ce dernier est un isomorphisme, donc il en est de même pour (1.3.1).

\textbf{Définition 1.4.} Soient $(X,A)$ un site commutativement localement annelé, $E$ un $A$-Module localement libre de type fini, $u:E\longrightarrow A$ un homomorphisme $A$-linéaire. On dit que $u$ est régulier si le complexe de Koszul $K_\bullet(u)$ défini par $u$ est acyclique en degrés $>0$. Soit $J$ un Idéal de $A$; on dit que $J$ est
